OpenFisca-France 176.1 (législation des revenus 2025), cas-types sans niche ; calculs de l'auteur (synthese/proposition\_ircsg.py, outils\_ircsg.py).\par Note : proposition : franchise de 8\,000 \euro{} par adulte sur les revenus d'activité, puis taux marginal de 15 \% au premier euro imposé, 15 \% à 20 000 \euro{}, 28,8 \% à 30 000 \euro{}, 50 \% à 231 000 \euro{}, en ligne droite entre ces points ; pensions exonérées jusqu'à 12 500 \euro{}, puis de moins en moins jusqu'à 25 000 \euro{} ; chaque point est le supplément d'IR, de CSG et de CRDS dû pour 1 000 \euro{} de revenu après cotisations en plus (salaire net de cotisations, ou pension brute), calculé tous les 100 \euro{} sans lissage. Le seul saut qui reste est un vrai effet de la règle : entre 12 500 et 25 000 \euro{} de pension, chaque euro de pension en plus retire un euro d'exonération, si bien que le retraité paie deux fois le taux du barème, et ce taux monte vite entre 20 000 et 25 000 \euro{}.
